% !PRINTED 142
\weberpage{142}
%<<CONT

\weq{6}{\begin{aligned}
\frac{H(v)}{\Theta(v)} &= \sqrt{\varkappa}\operatorname{sn}v,\\
\frac{H(v+K)}{\Theta(v)} &= \sqrt{\frac{\varkappa}{\varkappa'}}\operatorname{cn}v,\\
\frac{\Theta(v+K)}{\Theta(v)} &= \frac{1}{\sqrt{\varkappa'}}\operatorname{dn}v.
\end{aligned}}
The functions $\Theta$, $H$ have, as follows easily from (3), the
periodicity expressed by the following equations:
\weq{7}{\begin{aligned}
\Theta(v+2nK) &= \Theta(v),\\
\Theta(v+2niK') &= (-1)^n e^{\frac{n^2\pi K'}{K}-\frac{\pi inv}{K}}\Theta(v),
\end{aligned}}
\weq{8}{\begin{aligned}
H(v+2nK) &= (-1)^n H(v),\\
H(v+2niK') &= (-1)^n e^{\frac{n^2\pi K'}{K}-\frac{\pi inv}{K}}H(v).
\end{aligned}}
It is sometimes useful to combine the equations (3), (7), (8)
in the following way:
\weq{9}{\Phi(v+niK') = i^n e^{\frac{n^2\pi i}{4}\frac{K'}{K}-\frac{\pi inv}{2K}}\Psi(v),}
in which, if $n$ is \gesperrt{even}, $\Phi$ and $\Psi$ are both to be set equal to
$\Theta$ or both equal to $H$; if $n$ is odd, $\Phi=\Theta$, $\Psi=H$ or
$\Phi=H$, $\Psi=\Theta$.

From \S\,22 (2), (4) there result the addition formulae:
\weq{10}{\begin{aligned}
\Theta^2(0)\Theta(u+v)\Theta(u-v) &= \Theta^2(u)\Theta^2(v) - H^2(u)H^2(v),\\
\Theta^2(0)H(u+v)H(u-v) &= H^2(u)\Theta^2(v) - \Theta^2(u)H^2(v),
\end{aligned}}
for which one may also write:
\weq{11}{\begin{aligned}
\Theta^2(0)\Theta(u+v)\Theta(u-v) &= \Theta^2(u)\Theta^2(v)(1-\varkappa^2\operatorname{sn}^2u\operatorname{sn}^2v),\\
\Theta^2(0)H(u+v)H(u-v) &= \Theta^2(u)\Theta^2(v)\varkappa(\operatorname{sn}^2u-\operatorname{sn}^2v).
\end{aligned}}
We now proceed to carry over the theorems on the $\vartheta$-functions
to the elliptic functions, and we begin with
the addition theorem.

\begin{center}
\textbf{\S\,44. \gesperrt{Addition theorem of the elliptic functions}.}
\end{center}

From the four $\vartheta$-functions twelve quotients in all can be
formed, which by \S\,42 are expressible by elliptic functions.
If one forms these quotients for the argument
$u+v$, then by \S\,22 one can represent them by $\vartheta$-functions of $u$
and of $v$ separately, and hence also by the corresponding elliptic
 functions, and indeed always in such a way that the denominator
%CONT>>
